\documentclass{beamer}
\usepackage[slovene]{babel}
\usepackage[utf8]{inputenc}
\usepackage[T1]{fontenc}
\usepackage{lmodern}
\usepackage{amsmath, amssymb, bbm}

\usetheme{metropolis}
\beamertemplatenavigationsymbolsempty
\setbeamertemplate{caption}[numbered]

\usepackage{palatino}
\usefonttheme{serif}
\theoremstyle{definition}
\newtheorem{definicija}{Definicija}
% Krožne matrike
% Samo Primer
\title{Krožne matrike}
\author{Samo Primer}
\date{}
\institute[FMF]{Fakulteta za matematiko in fiziko}

% 
% Pozor: dokler ne dodate vsaj enega okolja za prosojnico, 
% se datoteka ne bo prevedla.
% 

\begin{document}
\begin{frame}
        \titlepage
\end{frame}


\begin{frame}
        \frametitle{Definicija}
% začetek definicije
\begin{definicija}
    Krožna matrika $n \times n$ je matrika oblike 
\begin{equation*}
  \left(
  \begin{matrix}
    c_0 & c_{n-1} & \cdots & c_2 & c_1 \\
    c_1 & c_0 & \cdots & c_{n-1} & c_2 \\
    c_2 & c_1 & \cdots & c_0 & c_{n-1} \\
    \vdots  & \vdots  & \ddots & \vdots & \vdots \\
    c_{n-1} & c_{n-2} & \cdots & c_1 & c_0
  \end{matrix}
  \right)
\end{equation*}
\end{definicija}       
% konec definicije
\end{frame}

\begin{frame}
        \frametitle{Lastnosti}
% začetek neoštevilčenega seznama
\begin{itemize}[<+->]
        % lastnost 1
        \item \label{lastnost 1}Lastni vektorji krožne matrike so 
        $v_j = (1, \omega_j, \omega_j^2, \ldots, \omega_j^{n-1})^T$, 
        kjer je $\omega_j=e^\frac{2\pi i j}{n}$ za $j=0, \ldots, n-1$.
        % lastnost 2
        \item \label{lastnost 2} Lastni vektorji krožne matrike so stolpci matrike enotske diskretne Fourierjeve transformacije, 
        ki jo lahko prikažemo kot Vandermondovo matriko.
        % lastnost 3
        \item \label{lastnost 3}Krožne matrike tvorijo komutativno algebro, ker je za poljubni dve matriki 
        $A$ in $B$, vsota $A + B$ tudi krožna matrika, prav tako je krožna matrika tudi njun produkt 
        $AB$, ter tudi velja $AB = BA$.
\end{itemize}
% konec neoštevilčenega seznama
\end{frame}
\begin{frame}
        \cite{gray}
\frametitle{Literatura}
\bibliographystyle{siam}
\bibliography{literatura}
\end{frame}



\end{document}

